\section{An explicit all-direction physical horizon and lattice budget}
\label{sec:eh-horizon}
We give a quantitative proof of finite horizon for the unchanged round
scatterers. It is independent of a simulated trajectory, a symbolic
model, a pressure theorem and the compactness argument retained in
Section~\ref{sec:uc-complexity}. It makes the finite geometric constants
in the physical source construction directly auditable.

\begin{lemma}[A row rectangle for a primitive lattice direction]
\label{lem:eh-rectangle}
Let $v$ be a primitive vector of the triangular lattice, $l=|v|$, and
$h=\sqrt3/(2l)$. For every $q\in\R^2$ there is a lattice center $z$
whose components relative to $q$, along $v/l$ and its unit normal,
satisfy $|x|\le l/2$ and $|y|\le h/2$. Thus for a line direction
making angle $\delta$ with $v$,
\begin{equation}\label{eq:eh-projections}
 |(z-q)\cdot e^\perp|
    \le \tfrac h2|\cos\delta|+\tfrac l2|\sin\delta|,
 \qquad |(z-q)\cdot e|\le\tfrac12\sqrt{l^2+h^2}.
\end{equation}
\end{lemma}
\begin{proof}
Primitivity permits $v$ to be completed to a lattice basis. The parallel
lattice rows are consequently spaced by the covolume $\sqrt3/2$ divided
by $l$. Choose a nearest row to $q$, and then a nearest lattice center
along that row. Its two coordinates obey the stated bounds, regardless
of the longitudinal offset of the row. Projection and Cauchy--Schwarz
give \eqref{eq:eh-projections}. All choices are in the original Euclidean
lattice; no linear map deforms the circular disks.
\end{proof}

\begin{theorem}[A uniform, explicit physical free-flight bound]
\label{thm:eh-horizon}
For every $R\in[9/20,47/100]$, every straight segment avoiding the
interiors of all radius-$R$ lattice disks has length at most
\begin{equation}\label{eq:eh-horizon}
 \tau_{\max}\le\sqrt{\frac{3847}{124}}<\frac{28}{5}<6.
\end{equation}
In particular the same bound holds for every genuine free flight and
all its grazing limits. The statement is simultaneous over all initial
points and all directions; no exceptional direction is discarded.
\end{theorem}
\begin{proof}
The dihedral symmetries of the triangular lattice reduce an unoriented
line direction to $0\le\vartheta\le\pi/6$. Put
$t=\tan\vartheta/\sqrt3\in[0,1/3]$. We use four primitive lattice
vectors, with the indicated ranges:
\[
\begin{array}{c|c|c}
\text{range of }t&v&|v|^2\\ \hline
{[0,1/60]}&(1,0)&1\\
{[1/60,1/12]}&(11/2,\sqrt3/2)&31\\
{[1/12,1/5]}&(5/2,\sqrt3/2)&7\\
{[1/5,1/3]}&(3/2,\sqrt3/2)&3
\end{array}
\]
The last three vectors have lattice coordinates $(5,1),(2,1),(1,1)$,
so they are primitive. For the horizontal vector the squared normal
width in \eqref{eq:eh-projections} is
\[
 W_0(t)^2=\frac3{16}\frac{(1+2t)^2}{1+3t^2}.
\]
This is increasing on $[0,1/60]$, since the logarithmic derivative of
$(1+2t)/\sqrt{1+3t^2}$ has numerator $2-3t>0$.
For $v=((2a+1)/2,\sqrt3/2)$, write $s=2a+1$ and
$N=a^2+a+1$. The normal width is
\[
 W_a(t)=\frac{\sqrt3}{4\sqrt{1+3t^2}}
       \left\{\frac{s+3t}{2N}+|st-1|\right\}.
\]
On each chosen interval for $a=5,2,1$ one has $t\le1/s$.
The expression in braces is positive and has derivative
$3/(2N)-s<0$. Since the denominator increases, $W_a$ is decreasing.
It therefore suffices to check its left endpoint. The four exact
squared widths and their positive gaps below $R_-^2=81/400$ are
\[
\begin{array}{c|c|c}
 v\text{ type}&\max W^2&81/400-\max W^2\\ \hline
\text{horizontal}&961/4804&1181/480400\\
a=5&13697401/73866304&31513139/1846657600\\
a=2&529/3136&2651/78400\\
a=1&75/448&393/11200
\end{array}
\]
Thus, for every direction and midpoint $q$, Lemma~\ref{lem:eh-rectangle}
provides a lattice center at \emph{strictly} less than $R_-$ normal
distance from the line. The longitudinal distance is at most
$\sqrt{l^2+3/(4l^2)}/2$. Among the four lengths its largest value is
$\sqrt{31+3/124}/2=\sqrt{3847/124}/2$.
If a free segment were longer than twice this number, the perpendicular
foot of the chosen center would lie inside the segment and inside an
open radius-$R_-$ disk, a contradiction. This proves the first
inequality in \eqref{eq:eh-horizon}. The comparison with $28/5$ follows
by squaring positive quantities:
$3847\cdot25=96175<97216=784\cdot124$.
Since the disk radius is only increased for $R\ge R_-$, the entire
argument is uniform over the original family. Strict normal clearance
in the covering argument includes limiting directions and grazing free
segments; no collision derivative at grazing is used.
\end{proof}

\begin{corollary}[Explicit enclosing candidate counts, not word counts]
\label{cor:eh-counts}
Keep the original constants $R_+=47/100$, $\delta_*=3/200$, the intrinsic
$\tau_+$ of the retained source papers, and every original event. Then
\[
 N_f\le150,\qquad N_\circ=N_*\le150,\qquad N_{\rm pair}\le162,
 \qquad J_f=2N_f+1\le301.
\]
Here $N_f$ is the first-hit candidate count of
\eqref{eq:uc-candidates}, $N_\circ$ is the one-mark/cross-flight count
of \eqref{eq:xf-candidates}, and $N_{\rm pair}$ is the difference count
in \eqref{eq:dc-constants}. In particular the angular first-hit partition
in Lemma~\ref{lem:ag-branches} has at most $1201$ intervals.
These are upper bounds for the intrinsic counts, not assertions that
all candidate centers can occur or that the smallest physical horizon
is $28/5$.
\end{corollary}
\begin{proof}
Theorem~\ref{thm:eh-horizon} allows the enclosing constant
$\overline\tau=28/5$. The corresponding three enclosing radii are
\[
 D_f^{\rm enc}=327/50,\qquad
 D_\circ^{\rm enc}=1311/200,\qquad
 D_{\rm pair}^{\rm enc}=657/100.
\]
Their squares are respectively $106929/2500$, $1718721/40000$ and
$431649/10000$. The first two lie in $[42,43)$ and the last in
$[43,44)$. A lattice point with integer coordinates $(a,b)$ has
squared length $Q(a,b)=a^2+ab+b^2$, which is integral. The enclosing
sets are therefore exactly $0<Q\le42$ and $0<Q\le43$.

For completeness this finite count does not depend on a floating-point
circle. The identity $Q=(a+b/2)^2+3b^2/4$ implies $|b|\le7$ when
$Q\le43$, and by symmetry $|a|\le7$. Direct solution of the quadratic
inequality gives the following inclusive integer ranges for $b\ge0$:
\[
\begin{array}{c|c|c|c|c}
 b& a\text{ for }Q\le42&\#&a\text{ for }Q\le43&\#\\ \hline
0&[-6,6]&13&[-6,6]&13\\
1&[-6,5]&12&[-7,6]&14\\
2&[-7,5]&13&[-7,5]&13\\
3&[-7,4]&12&[-7,4]&12\\
4&[-7,3]&11&[-7,3]&11\\
5&[-7,2]&10&[-7,2]&10\\
6&[-6,0]&7&[-7,1]&9\\
7&[-5,-2]&4&[-6,-1]&6
\end{array}
\]
Negative rows have the same cardinalities under $(a,b)\mapsto(-a,-b)$.
The totals, excluding the origin, are
$12+2(12+13+12+11+10+7+4)=150$ and
$12+2(14+13+12+11+10+9+6)=162$.
The intrinsic candidate sets are subsets of these enclosing sets,
which proves all three inequalities. The stated branch bounds now
follow from the already proved formulas $2N_f+1$ and $8N_*+1$.
No collision word, pairing sheet, homogeneous grazing strip or
probability normalization is counted by these finite numbers.
\end{proof}

This quantitative result replaces no dynamical hypothesis by a finite
simulation. It supplies the geometric finite-horizon input used in the
retained source and curve-growth proofs. It does not supply the
long-time marked pressure, the sum at a common physical roof, or the
collision-normalized positive-source estimate.
